\documentclass[11pt,a4paper]{article}

\usepackage[a4paper,margin=2.5cm]{geometry}
\usepackage{amsmath}
\usepackage{amssymb}
\usepackage{graphicx}
\usepackage{hyperref}

\title{
Unsigned Natural-Number Wavefront Propagation
and Emergent Euclidean Geometry
}

\author{
Alexandre Furtado Neto\\
Department of Physics (Alumnus), UNESP, BR\\
\texttt{alexandre.com@yahoo.com}
}

\date{\today}
\begin{document}

\maketitle

\begin{abstract}
We present a minimal cellular automaton on a cubic lattice in which Euclidean quadratic geometry emerges from purely local updates using only unsigned integer arithmetic. Each cell stores a single scalar field \(r^2\), representing the accumulated squared distance from the origin. Propagation is governed by the discrete quadratic increment identity \((x+1)^2 - x^2 = 2x + 1\) (and cyclic permutations), applied locally along Cartesian directions.

The system operates entirely over the natural numbers \(\mathbb{N}\), requires no floating-point operations, signed coordinates, or global computations, and exactly reconstructs the Euclidean metric \(r^2 = x^2 + y^2 + z^2\). The resulting wavefronts exhibit stable radial expansion and contraction, suggesting applications in discrete causal models, arithmetic geometry, and massively parallel hardware implementations.
\end{abstract}

\section{Introduction}
Conventional computational models of geometry typically rely on continuous coordinates, signed arithmetic, floating-point operations, or explicit global metric evaluation \cite{danielsson1980,sethian1996}. In contrast, this work demonstrates that Euclidean structure can emerge from strictly local propagation rules defined solely over natural numbers.

The proposed model propagates a scalar field across a cubic lattice using only nearest-neighbor interactions and monotonic unsigned additions. Remarkably, the global field converges exactly to the Euclidean quadratic metric without ever computing the expression \(x^2 + y^2 + z^2\) directly.

\section{Discrete Quadratic Increment}

The foundation of the model is the elementary finite-difference identity:
\[
(x+1)^2 - x^2 = 2x + 1,
\]
with analogous expressions for the \(y\) and \(z\) directions. This identity enables exact incremental reconstruction of the quadratic form through local updates alone.

\section{Lattice Model and Propagation Rule}

Consider a cubic lattice of size \(L \times L \times L\) centered at \((MID, MID, MID)\). Each cell stores a single unsigned integer \(r^2\), initialized to \(\infty\) except at the origin, where \(r^2 = 0\).

Absolute distances are obtained locally without signed arithmetic:
\[
|x| = 
\begin{cases}
x - MID & \text{if } x > MID, \\
MID - x & \text{otherwise}.
\end{cases}
\]

For each cell, the algorithm attempts propagation to its six Cartesian neighbors. The incremental cost along the \(x\)-direction, for example, is \(\Delta_x = 2|x| + 1\). The propagated value
\[
r^2_{\text{new}} = r^2_{\text{current}} + \Delta
\]
updates the neighbor if smaller. The process is strictly monotonic and uses only additions and comparisons.

\section{Emergent Euclidean Geometry}

Although the propagation rule never evaluates the global expression \(r^2 = x^2 + y^2 + z^2\), the scalar field converges exactly to this quadratic metric at every lattice point. Numerical validation on a \(L=121\) lattice (\(R=60\)) confirms zero maximum and RMS errors in both quadratic and radial distances across 904,089 interior points. This exact reconstruction emerges solely from local unsigned finite-difference accumulation.

\section{Emergent Causal Fronts}

While the underlying field is quadratic, temporal selection of isosurfaces \(r^2 = f(t)\) with linearly increasing \(f(t)\) produces an observable shell that expands and contracts with approximately constant radial velocity (\(r(t) \propto t\)). This behavior arises without explicit velocity fields, momentum, or differential equations, suggesting the emergence of discrete causal horizons from purely arithmetic propagation.

\section{Relation to Existing Work}

The approach shares conceptual ground with Euclidean distance transforms \cite{danielsson1980}, fast marching methods \cite{sethian1996}, and Dijkstra-style propagation \cite{dijkstra1959}. In particular, Case et al. \cite{case2001lattice} demonstrated expanding wavefronts on cubic lattices that approximate Euclidean space. However, their model relies on multiplication operations within the cell update rule, thereby increasing the complexity of the underlying finite-state machine. In contrast, the present model achieves exact quadratic metric reconstruction using only additions and comparisons, while additionally featuring pulsating (expanding and contracting) behavior through a deterministic external threshold.

\section{Code Availability}

The complete implementation, including real-time visualization and precision analysis, is available as open-source software \cite{neto2026arithmetico}:

\begin{itemize}
    \item Repository: \url{https://github.com/automaton3d/arithmetico}
    \item License: MIT
\end{itemize}

\section{Discussion and Outlook}

This unsigned natural-number formulation suggests that Euclidean geometry may be reconstructed dynamically from local causal arithmetic rather than presupposed as a background structure. The model offers a primitive framework for digital physics, causal set theory, and emergent metric theories.

Future directions include anisotropic and non-Euclidean metrics, curvature, higher dimensions, dynamic perturbations, and hardware implementations.

\section{Conclusion}

We have introduced a minimal unsigned wavefront propagation algorithm that exactly reconstructs Euclidean quadratic geometry from local arithmetic updates over natural numbers. The emergence of both exact metric structure and stable causal fronts from such elementary rules highlights the power of discrete arithmetic propagation as a foundation for geometry and causal modeling.

\begin{thebibliography}{9}
\bibitem{danielsson1980} Per-Erik Danielsson. \textit{Euclidean Distance Mapping}. Computer Graphics and Image Processing, 1980.
\bibitem{sethian1996} J. A. Sethian. \textit{A Fast Marching Level Set Method for Monotonically Advancing Fronts}. Proc. Natl. Acad. Sci., 1996.
\bibitem{dijkstra1959} Edsger W. Dijkstra. \textit{A Note on Two Problems in Connexion with Graphs}. Numer. Math., 1959.
\bibitem{case2001lattice} J. Case, D. S. Rajan, and A. M. Shende. Lattice computers for approximating euclidean space. \textit{Journal of the ACM}, 48(1):110--144, 2001.
\bibitem{neto2026arithmetico} A. Neto. \textit{Arithmetico: Euclidean Geometry from Unsigned Integer Arithmetic}. GitHub repository, 2026. \url{https://github.com/automaton3d/arithmetico}
\end{thebibliography}

\end{document}
